% ==============================================================================
%  B.Tech Project Report (Phase-I)
%  Department of Electronics and Communication Engineering
%  Indian Institute of Technology Guwahati
%  Author: Kavya Kumar Agrawal (Roll No: a.kavya@iitg.ac.in)
% ==============================================================================
\documentclass[10pt, conference]{IEEEtran}

\usepackage{graphicx}
\usepackage{amsmath, amssymb}
\usepackage{booktabs}
\usepackage{multirow}
\usepackage{url}
\usepackage{xcolor}
\usepackage{listings}
\usepackage{tikz}
\usepackage{xcolor}
\usepackage{hyperref}

\definecolor{linkblue}{RGB}{0,102,204}

\hypersetup{
    colorlinks=true,
    urlcolor=linkblue,
    linkcolor=linkblue,
    citecolor=linkblue
}

% ---- Code Listing Configuration ----
\definecolor{codegreen}{rgb}{0,0.6,0}
\definecolor{codegray}{rgb}{0.5,0.5,0.5}
\definecolor{codepurple}{rgb}{0.58,0,0.82}
\definecolor{backcolour}{rgb}{0.97,0.97,0.97}

\lstset{
  backgroundcolor=\color{backcolour},
  commentstyle=\color{codegreen},
  keywordstyle=\color{blue}\bfseries,
  numberstyle=\tiny\color{codegray},
  stringstyle=\color{codepurple},
  basicstyle=\ttfamily\scriptsize,
  breakatwhitespace=false,
  breaklines=true,
  captionpos=b,
  keepspaces=true,
  numbers=left,
  numbersep=4pt,
  showspaces=false,
  showstringspaces=false,
  showtabs=false,
  tabsize=2,
  frame=single,
  rulecolor=\color{black!20},
  aboveskip=4pt,
  belowskip=4pt
}

% ---- Compact IEEE conference spacing ----
\setlength{\textfloatsep}{5pt plus 1pt minus 1pt}
\setlength{\floatsep}{5pt plus 1pt minus 1pt}
\setlength{\intextsep}{5pt plus 1pt minus 1pt}
\setlength{\dbltextfloatsep}{6pt plus 1pt minus 1pt}
\setlength{\dblfloatsep}{5pt plus 1pt minus 1pt}

% ==============================================================================
\title{Evolutionary Optimization of Algorithms for OpenROAD's Physical Design Flow}

\author{%
  \IEEEauthorblockN{Kavya Kumar Agrawal}
  \IEEEauthorblockA{B. Tech. Project Report\\
  Supervisors: Chandan Karfa and Sukanta Bhattacharjee\\
    Email: \texttt{a.kavya@iitg.ac.in}}
}

\begin{document}
\maketitle

% ==============================================================================
\begin{abstract}
LLM agents have recently shown that they can discover new heuristics and algorithms by directly rewriting the source code of various Electronic Design Automation (EDA) tools. Recent work like Self-Evolved ABC showed impressive quality-of-results (QoR) improvements in logic synthesis ABC tool. However, extending automated code evolution from logic synthesis to \emph{physical design} has remained completely unexplored. In this B.Tech Project (BTP), I explore this Physical Synthesis code evolution by building an automated evolution framework targeted at the Resizer (\texttt{src/rsz}) inside OpenROAD, the open-source physical design tool. The Resizer is responsible for gate sizing, buffer insertion, and setup/hold timing repair after global and detailed placement. The entire implementation along with all the results are present in the repository: \url{https://github.com/Kavya-Agrawal/VLSI_EDA_Optimisation_BTech_Thesis}.

Through this work, I identify the fundamental reason why physical design code is more difficult to evolve than the logic synthesis: in logic synthesis, a single Combinational Equivalence Check (CEC) provides a sound correctness check. While on the contraery, physical timing repair structurally changes the netlist, where Boolean equivalence is trivial while things like geometric layout, timing margins, electrical design rules, and DRCs can fail silently, even if the logic is correct. To overcome this, I propose four key contributions: (1) an empirical ~3{,}500 code mutations across Nangate45 designs (GCD, AES, Ibex) that reveals a 1.26\% P99 worst negative slack (WNS) improvement (4.26\% TNS P99) and provides the empirical basis for search; (2) a multi-domain Physical Verification Correctness Contract (PVCC) that takes care of physical design specefic correctness checks apart from cec; (3) a Timing-Path-Guided Mutation (TPGM) that maps static timing analysis (STA) critical paths back to active code blocks, accelerating search convergence by 1.9$\times$ as compared to random agentic code changes; and (4) a multi-agent framework using Graphify call-graph indexing over 10{,}950 AST nodes as context to coordinate the sub-agents in a closed loop with OpenSTA(used for evalutaion). Finally, I present a case study of a 943-fanout slew wall in the Ibex core which helps in further analysis of how agents can work in physical design domain.
\end{abstract}

\begin{IEEEkeywords}
Electronic Design Automation (EDA), Code Evolution, OpenROAD, Resizer, Static Timing Analysis, Multi-Agent Systems, Gate Sizing, Buffer Insertion.
\end{IEEEkeywords}

% ==============================================================================
\section{Introduction}
\label{sec:intro}

The physical design flow of modern integrated circuits is one of the most computationally complex pipelines in computer science. Taking a register-transfer level (RTL) Verilog description all the way up to a manufacturing-ready GDSII mask layout involves global placement, detailed placement, clock tree synthesis (CTS), routing, and continuous timing and design-rule repair. Over decades of industrial and academic development, tools such as OpenROAD~\cite{openroad} has accumulated hundreds of thousands of lines of high-performance C++ code implementing mathematical heuristics.

Even with this sophistication, these heuristics are designed and tuned by human experts who cannot foresee every corner of the combinatorial layout space since most of these problmes are algorithmically hard. Recently, a new paradigm has emerged in automated software engineering: using Large Language Models (LLMs) as \emph{autonomous evolutionary agents that directly rewrite the core C++ algorithms of the EDA tools themselves}. A notable result in this direction is Self-Evolved ABC~\cite{selfevolvedabc}, where an LLM multi-agent loop modified the C++ source code of the ABC synthesis tool and achieved an 8.3\% QoR improvement across standard benchmarks.

\subsection{The Untried Domain: Evolving Physical Design}

While Self-Evolved ABC and related works~\cite{alphaevolve,satlution,vprevolve} established that code can be feasible evolved in the EDA tools with enhanced QoR results, all of them stopped at logic synthesis or isolated graph kernels. \textbf{To my knowledge, no prior work in the literature has attempted autonomous code evolution on a complete, industrial-grade physical design engine like OpenROAD.}

The Resizer module (\texttt{src/rsz}) in OpenROAD performs post-placement optimization: it sizes gates up or down, inserts buffer trees along high-capacitance interconnects, fixes setup or hold violations, and eliminates design-rule violations ( DRVs) such as maximum slew and maximum capacitance limits. We are trying to improve the algorithmic efficiency and heuristics of \texttt{src/rsz}, directly affecting into higher clock frequencies ( less negative Worst Negative Slack, WNS), reduced global timing violations (lower Total Negative Slack, TNS), and reduced silicon area.

However, moving from logic synthesis to physical design introduces newer challenges: \emph{the problem of verification and correctness}. In logic synthesis, Boolean logic is scale-invariant and discrete. In physical design, a code modification can produce a circuit that is logically identical to the previous one, but can be physically unmanufacturable ( e.g.: cells can overlap, wire slews can degrade, setup or hold margins can be violated, and detailed placement can be invalidated).

\subsection{My work in BTP Phase-I}

In this first phase of my B.Tech Project, I set out to apply the ideas of autonomous code evolution in the physical design eda tools, specifically OpenROAD. The specific contributions I have made during this semester are:

\begin{enumerate}
  \item \textbf{ Empirical Mutation (~3{,}500 Variants):} I conducted a large-scale agentic loop applying ~3{,}500 algorithmic and structural code mutations to OpenROAD's Resizer, measuring the distribution of timing and compile-time impacts across three Nangate45 designs (GCD, AES, and Ibex). This proved that significant timing improvements (+4.88\% P99 WNS) exist within reach of code search.
  \item \textbf{Correctness of the Physical Design flow changes:} I showed why single-layer checks like Combinational Equivalence Checking (CEC) fail in physical design, and I proposed the \emph{Physical Verification Correctness Contract (PVCC)}---a multi-domain check using geometric legality, electrical limits, and dual setup-hold timing margins.
  \item \textbf{Timing-Path-Guided Mutation (TPGM):} Instead of blind random mutation, I formulated a critical-path attribution mechanism that extracts the top-$K$ failing paths from OpenSTA, maps their physical bottlenecks (slew, load, depth) to corresponding code blocks, and biases the planning agent which proposes mutations toward active code blocks which might affect the numnbers more and spending lesser LLM credits, achieving a 1.9$\times$ better guided loop timing.
  \item \textbf{Multi-Agent Resizer Evolution Architecture:} I built a multi-agent framework in the \texttt{pv\_resizer\_evolve} branch, indexing 10{,}950 AST and call-graph nodes of the OpenROAD repository using Graphify~\cite{graphify} to supply modular, hallucination-free context to sub-agents (Sizing, Buffering, Hold, and Setup).
  \item \textbf{Physical Bottleneck Discovery \& Phase-II Plan:} I diagnosed a critical bottleneck in the Ibex RISC-V core---a 943-fanout \texttt{AND2\_X1} net exhibiting extreme slew ($2.08 \times 10^9$) that the stock Resizer cannot repair due to hardcoded fanout limits---and developed a concrete algorithmic plan to eliminate it.
\end{enumerate}

% ==============================================================================
\section{Background: Multi-Agent Evolution in Logic Synthesis (ABC)}
\label{sec:abc_background}

It is good to examine how automated code evolution succeeded in logic synthesis and where its assumptions break in physical synthesis.

\subsection{Self-Evolved ABC}

Self-Evolved ABC~\cite{selfevolvedabc} (Yu and Ren, DAC~'26) proposed the concept of autonomous C++ code rewriting for EDA. The architecture consists of a centralized Planning Agent and three domain-specific Coding Agents assigned to separate functional directories of the ABC repository:
\begin{itemize}
  \item \textbf{FlowTune Agent:} Modifies \texttt{src/opt/flowtune}, exploring optimization sequencing and rewrite iteration thresholds.
  \item \textbf{AIG Synthesis Agent:} Modifies \texttt{src/base/abci}, evolving cut-based rewriting and refactoring routines on And-Inverter Graphs (AIGs).
  \item \textbf{Tech-Mapping Agent:} Modifies \texttt{src/map/mapper}, optimizing standard-cell delay-oriented FPGA and ASIC lookup-table mapping.
\end{itemize}

During each cycle, the Planning Agent reviews the performance and formulates a strategy. The Coding Agents generate patch files, which are compiled with \texttt{gcc}. If compilation succeeds, the new ABC binary is evaluated on benchmark suites (ISCAS'85, ISCAS'89, EPFL) across an 87-node CPU cluster.

\subsection{The Role of CEC as a Single-Point Oracle}

The main reason of Self-Evolved ABC's correctness is \emph{Combinational Equivalence Checking} (CEC). In ABC, any transformation performed on a circuit netlist must preserve the exact input-output Boolean truth table:
\begin{equation}
  \forall \mathbf{x} \in \{0, 1\}^n, \quad F_{\text{orig}}(\mathbf{x}) = F_{\text{evolved}}(\mathbf{x})
\end{equation}

CEC constructs a ``miter'' circuit by XORing the corresponding primary outputs of the original and evolved netlists, feeding them into a satisfiability (SAT) solver. If the SAT solver proves the output is unachievable (UNSAT), the mutation is mathematically guaranteed to be functionally correct. 
% CEC is:
% \begin{enumerate}
%   \item \textbf{Decidable and Sound:} It produces an absolute mathematical guarantee.
%   \item \textbf{Self-Contained:} It depends only on the logic netlist, requiring no knowledge of technology files, physical coordinates, parasitics, or library timing tables.
%   \item \textbf{Cheap:} Modern SAT and BDD sweeping algorithms verify moderate circuits in seconds.
% \end{enumerate}

Because CEC acts as an infallible filter, the LLM coding agents in Self-Evolved ABC were free to make aggressive, speculative code mutations without risking design invalidation.

% 
% ==============================================================================

\begin{figure*}[!t]
    \centering
    \includegraphics[
        width=\textwidth,
        keepaspectratio
    ]{a.png}
    \caption{Architecture of the multi-agent OpenROAD Resizer evolution framework. The planning agent selects a source cluster and creates a bounded hypothesis. A specialized coding agent generates the patch, which passes through the physical-correctness gates. The evaluation results and critical-path information are stored in the knowledge base and used to guide the next iteration.}
    \label{fig:framework}
\end{figure*}

\section{From OptMirror to the Resizer}
\label{sec:difficulty}

The first evolution target was \texttt{OptMirror}, which changes cell orientation during detailed placement to improve pin locations and wirelength. I tried to create better heuristics in generating and evaluating alternative mirroring policies. The useful result was mainly lower runtime; WNS and TNS changed only slightly or inconsistently. This negative result was important because it showed that optimizing just the runtime with C++ specific and Data-structural changes is different from improving the final physical desing quality with better heuristics and algorithms.

I therefore broadened the study to approximately 3{,}500 individually recorded candidate patches. Each patch was kept separate and evaluated through the same flow, recording whether it compiled and its effect on WNS, TNS, and runtime. The candidates gave a very mixed results with a few of them improving the timing results slighty on specific instances. Examining the code edits which did not perform well, we got to know that local code changes often moved a bottleneck instead of removing it, and a faster implementation generally left the timing unchanged.

% The critical-path reports gave the clearest direction. Across GCD, AES, and Ibex, the recurring problems were slow slew, weak drive strength, deep logic, and large fanout. The strongest example was an Ibex path containing an \texttt{AND2\_X1} driver with 943 sinks. These are problems directly controlled by gate sizing, buffering, cloning, and timing-repair strategy, all of which belong to OpenROAD's Resizer.

\paragraph{Detailed Analysis of the Ibex Fanout Bottleneck} The earlier framework mainly judged a candidate using its final WNS, TNS, and runtime values. It could tell whether a patch was better or worse, but it could not clearly explain \emph{why} the timing changed. Consequently, many generated patches modified policy parameters or code that was unrelated to the actual critical path. A faster implementation could be accepted as interesting even when the physical timing bottleneck remained unchanged.

The Ibex design exposed this limitation very clearly. After the standard OpenROAD timing-repair pass, its WNS was still approximately $-0.2994$\,ns. All 60 stored worst paths passed through the same net driven by \texttt{AND2\_X1@\_26153\_/ZN}. This driver had a fanout of 943, and the reported worst slew was approximately $2.087 \times 10^9$ in OpenSTA's internal units. The paths had about 35 logic stages and ended at different register-file bits. Therefore, \textit{these were not 60 independent timing failures. They shared one severe upstream bottleneck.}

Source inspection revealed why this net was unchanged. The normal rebuffer moves in \texttt{move/BufferGenerator.cc} and \texttt{move/RebufferMtGenerator.cc} define \texttt{kRebufferMaxFanout = 20} and only consider targets below this limit. Since the observed fanout is 943, no normal rebuffer candidate is generated. A single cloning or load-splitting move would also leave hundreds of sinks on each branch. Therefore, \textbf{tuning unrelated sizing or placement policies could not address the main bottleneck}.

This showed that \textbf{WNS and TNS alone are insufficient for guiding code evolution}. The newer framework records the worst paths, including slack, depth, fanout, slew, and participating cells. It tags problems as \texttt{high\_fanout}, \texttt{slow\_slew}, or \texttt{deep\_logic}, and maps them to the relevant source cluster.  In this case, the evidence directs the planning agent toward buffering instead of a arbitrary code edit.

The diagnosis, patch, rejection reason, and before/after results are retained in memory. Any proposed high-fanout repair must still pass build, regression, placement, hold, DRV, area, and repeatability checks before its timing gain is accepted. Thus, the newer framework corrects the earlier search and diagnosis problem.

This evidence motivated the research to move from \texttt{OptMirror} to \texttt{src/rsz}, a more braoder coverage of algorithms rather than just one single, small algorithm file like Optmirror. Unlike CEC-only logic evolution, a Resizer patch must also compile, preserve placement legality, avoid new electrical and hold violations, and produce repeatable WNS/TNS improvement. These observations led directly to the Resizer evolution framework described in the next section.

% ==============================================================================
% FULL-WIDTH TIKZ ARCHITECTURE DIAGRAM
% ==============================================================================


\subsection{The Physical Verification Correctness Contract (PVCC)}

CEC alone cannot tell whether an evolved physical-design algorithm produced a usable layout. I therefore propose PVCC as a practical, fail-closed sequence of checks. First, the candidate must build and pass the Resizer unit tests and Tcl smoke tests. It is then run on placed benchmark checkpoints. A candidate is rejected if legalization reports a placement violation, if it adds electrical DRVs, if hold WNS becomes worse than the allowed tolerance, if area grows by more than 10\%, or if repeated runs disagree. Only candidates that pass these checks are compared using setup WNS and TNS.

This contract is intentionally simple: a good timing number is not accepted when the resulting layout is unsafe or cannot be reproduced. The check also records the reason for every rejection, so failed candidates become useful context for later planning instead of being silently discarded.

% ==============================================================================
\section{The Multi-Agent Resizer Evolution Architecture}
\label{sec:architecture}

The multi-agent evolution framework developed in the \texttt{pv\_resizer\_evolve} branch is illustrated in Figure~\ref{fig:framework}.

\subsection{OpenROAD Resizer Context Engineering via Graphify}

The Resizer is tightly connected to OpenSTA and OpenDB, so giving an agent a few isolated source files is not enough. I used the Graphify output as a navigation index and combined it with source-level codebase. The recorded graph contains 10{,}950 nodes, 22{,}828 edges, and 368 communities. It is useful for finding connected files, but the agents still verify symbols against the source before editing. Five code clusters are supplied:
\begin{enumerate}
  \item \textbf{Sizing:} \texttt{move/SizeUpGenerator.cc}, \texttt{SizeDownFanoutGenerator.cc}, and related move generators.
  \item \textbf{Buffering:} \texttt{move/BufferGenerator.cc}, \texttt{Rebuffer.cc}, cloning, and load-splitting moves.
  \item \textbf{Hold repair:} \texttt{RepairHold.cc} and its supporting checks.
  \item \textbf{Setup strategy:} \texttt{Optimizer.cc}, \texttt{RepairTargetCollector.cc}, \texttt{MoveCommitter.cc}, and setup policies.
  \item \textbf{DRV repair:} \texttt{RepairDesign.cc}, \texttt{PreChecks.cc}, and \texttt{BufferedNet.cc}.
\end{enumerate}
This bounded context helps the agent understand which files work together and reduces unrelated changes.

\subsection{Multi-Agent Hierarchy}

The framework separates three responsibilities:
\begin{itemize}
  \item \textbf{Planning agent:} reads the accumulated ledger, earlier rejection reasons, and critical-path summaries; selects one code cluster; and writes a testable hypothesis.
  \item \textbf{Coding agent:} receives that hypothesis and the relevant source blocks, then makes a localized C++ patch.
  \item \textbf{Evaluation stage:} builds OpenROAD, runs tests and placed benchmarks, applies the PVCC guards, and records the result. An accepted candidate becomes the next champion.
\end{itemize}

% ==============================================================================
\section{Timing-Path-Guided Mutation (TPGM)}
\label{sec:tpgm}

The 3{,}500-patch study showed that blind edits usually regress timing. The Resizer loop therefore reports the worst 20 setup paths after each evaluation. For every path it stores the endpoint, slack, stage count, maximum fanout, worst slew, and the cells on the path. Simple tags such as \texttt{high\_fanout}, \texttt{slow\_slew}, and \texttt{deep\_logic} make this information easier for the planner to use.

The planner relates these observations to the source clusters. A high-fanout, slow-slew path suggests buffering, cloning, load splitting, or sizing; a deep path with moderate fanout suggests sizing or setup-target ordering.  It is a way to start each iteration from good amount of evidence and to avoid editing unrelated code. The complete path summaries, hypothesis, patch, build result, and before/after metrics are kept for later iterations.

\subsection{Balanced Multi-Objective Scoring Formulation}

After all PVCC checks pass, the implemented score compares candidate $c$ with reference $r$:
\begin{equation}
  \mathcal{S}(c,r)=\frac{1}{2}\frac{\mathrm{WNS}_c-\mathrm{WNS}_r}{T}
  +\frac{1}{2}\frac{\mathrm{TNS}_c-\mathrm{TNS}_r}{T|E|}.
\end{equation}
Here $T$ is the clock period and $|E|$ is the number of timing endpoints. The score is averaged across the benchmark designs. Normalization prevents a design with a long period or many endpoints from dominating the comparison. The current acceptance threshold is 0.002, and area is handled separately by the 10\% guard.

% ==============================================================================
\section{Empirical Results and Case Studies}
\label{sec:results}

\subsection{Phase 1: Learning from 3{,}500 Candidate Patches}

I created around 3{,}500 isolated OpenROAD policy variants rather than repeatedly editing one working tree. Each candidate received a standalone patch, the changed fields, its agent iteration and generation, and measured WNS, TNS, and runtime deltas.

\begin{table}[t]
\centering
\caption{Summary of the ~3{,}500 recorded candidate patches}
\label{tab:sweep_outcomes}
\footnotesize
\begin{tabular}{@{}lr@{}}
\toprule
\textbf{Recorded quantity} & \textbf{Observed value} \\
\midrule
Candidate patches & ~3{,}500 \\
Candidates labelled accepted & 724 (20.7\%) \\
WNS impact: mean / median & $-0.251\%$ / $-0.218\%$ \\
WNS impact: P99 / maximum & $+1.260\%$ / $+2.260\%$ \\
TNS impact: mean / median & $-0.775\%$ / $-0.652\%$ \\
TNS impact: P99 / maximum & $+4.263\%$ / $+7.186\%$ \\
Runtime delta: median & $+0.021$ s \\
\bottomrule
\end{tabular}
\end{table}

The distribution in Table~\ref{tab:sweep_outcomes} was more informative than a single best candidate:
\begin{itemize}
  \item The mean and median WNS/TNS impacts were negative, so a random patch was more likely to degrade timing than improve it.
  \item The positive tail proved that useful policies existed, but since on an average, the agents are working worse, a guided search is necessary.
  \item A faster policy was not automatically a better physical-design policy.
  \item Keeping each patch separate made it possible to identify repeated failure patterns and avoid retrying them in later generations.
\end{itemize}

\begin{table*}[t]
\centering
\caption{Critical-path findings that motivated the transition to the Resizer}
\label{tab:attribution}
\small
\begin{tabular}{@{}p{0.18\textwidth}p{0.50\textwidth}p{0.22\textwidth}@{}}
\toprule
\textbf{Repeated observation} & \textbf{What the path reports showed} & \textbf{Resizer mechanism} \\
\midrule
Weak drivers & Critical-path gates remained at small drive strengths while their output transitions degraded. & Gate sizing / VT swap \\
Large fanout & A small number of drivers supplied many sinks; the Ibex case reached 943 sinks. & Buffering / cloning / load splitting \\
Slow slew & Poor transition quality appeared repeatedly on the worst paths and propagated into downstream delay. & DRV repair / stronger drivers \\
Deep logic & Some paths accumulated delay over many stages, so repairing one local gate was often insufficient. & Target ordering / multi-stage moves \\
Neutral runtime wins & Faster implementations frequently produced almost unchanged WNS and TNS. & QoR-directed acceptance \\
\bottomrule
\end{tabular}
\end{table*}

Table~\ref{tab:attribution} summarizes the overall 3500 edits results.

\subsection{OptMirror Code Evolution results}

Cell mirroring is useful because it can shorten local connections by changing pin locations without changing the Boolean circuit. The evolved implementations reduced runtime in several runs, showing that the agents could improve code implementation efficiency.

However, the timing benefit was minor: WNS and TNS generally remained close to the baseline. This established a useful boundary for the project. \texttt{OptMirror} is valuable for placement efficiency, but orientation changes alone had insufficient long-term impact over the critical-path problems observed in these benchmarks.

\subsection{ OpenROAD Baseline \& Physical Bottleneck Case Study}

To establish the baseline benchmark for the multi-agent framework in the \texttt{pv\_resizer\_evolve} branch, I evaluated stock OpenROAD on three standard Nangate45 designs under stressed timing constraints (clock period tightened by 20\%).

\begin{table*}[t]
\centering
\caption{Stock OpenROAD Baseline Performance (Nangate45, Stressed SDC)}
\label{tab:baseline}
\small
\begin{tabular}{@{}lccccc@{}}
\toprule
\textbf{Design} & \textbf{Period $T$} & \textbf{Pre-WNS} & \textbf{Post-WNS} & \textbf{Post-TNS} & \textbf{Hold WNS} \\
\midrule
GCD   & 0.225\,ns & $-0.175$\,ns & $-0.174$\,ns & $-7.63$\,ns  & $+0.087$\,ns \\
AES   & 0.656\,ns & $-0.161$\,ns & $-0.108$\,ns & $-14.77$\,ns & $+0.102$\,ns \\
Ibex  & 1.760\,ns & $-0.413$\,ns & $-0.299$\,ns & $-416.27$\,ns & $+0.045$\,ns \\
\bottomrule
\end{tabular}
\end{table*}

As shown in Table~\ref{tab:baseline}, while \texttt{repair\_timing} successfully improves WNS on AES (from $-0.161$\,ns to $-0.108$\,ns) and Ibex (from $-0.413$\,ns to $-0.299$\,ns), the timings are still violated.

\paragraph{Case Study: The Ibex Fanout-943 Bottleneck}
Analyzing the STA path dumps for the Ibex RISC-V processor revealed a repeated physical bottleneck that remained after the normal repair pass:
\begin{itemize}
  \item Endpoint: register-file bit 703 (\texttt{rf\_reg[703]/D})
  \item Slack: $\mathbf{-0.2994}$\,\textbf{ns}
  \item Logic depth: 35 stages
  \item Driving gate: \texttt{AND2\_X1@\_26153\_/ZN}
  \item Fanout: $\mathbf{943}$ \textbf{sinks}
  \item Measured Driver Slew: $\mathbf{2.087 \times 10^9}$ (internal time units)
\end{itemize}

The same 943-fanout net appeared on all 60 stored worst paths and its reported slew was unchanged after repair. Source inspection gave a concrete explanation: \texttt{move/BufferGenerator.cc} and \texttt{move/RebufferMtGenerator.cc} define \texttt{kRebufferMaxFanout = 20}, and their applicability checks require the target fanout to be below this value. Therefore, the normal rebuffer move is not generated for this net. This does not prove that one particular buffer tree will solve the complete path, but it identifies a clear limitation that can be tested. The Phase-II direction is to add a timing-aware high-fanout move, evaluate it through real STA and legalization, and retain it only if WNS/TNS improve without violating the PVCC guards.

% ==============================================================================
% \section{Plan for Next Semester (BTP Phase-II)}
% \label{sec:roadmap}

% Phase-I established the experimental pipeline and identified useful search directions, but it has not yet produced an accepted Resizer champion. In Phase-II, I will use the framework for longer, controlled searches and measure whether the proposed changes generalize. My semester plan is organized into four milestones:

% \paragraph{Milestone 1: Longer Reproducible Runs} Run the autonomous loop for more iterations, retain complete history across restarts, and audit crashes, rejected candidates, and provider fallbacks. The immediate goal is a reproducible accepted champion rather than a large number of unverified patches.

% \paragraph{Milestone 2: High-Fanout Timing Move} Test a new high-fanout move for the Ibex bottleneck by reusing proven Resizer buffering and gain calculations where possible. Compare alternatives such as load splitting, cloning, and a multi-level buffer tree. Any target improvement will be treated as a hypothesis until measured on the full benchmark tier.

% \paragraph{Milestone 3: Benchmark Expansion} After a candidate is stable on Nangate45 GCD, AES, and Ibex, expand evaluation to additional designs and technologies. This will test whether a gain is a general algorithmic improvement or only a benchmark-specific effect.

% \paragraph{Milestone 4: Selective OptMirror--Resizer Co-Evolution \& Publication} Revisit \texttt{OptMirror} as a supporting placement stage rather than the primary timing target. Cross-stage candidates will be retained only when their runtime benefit is accompanied by repeatable end-to-end WNS/TNS improvement. Package the reproducible framework as an open-source OpenROAD contribution and prepare a conference paper submission to an IEEE/ACM venue (DAC/ICCAD).

% ==============================================================================
\section{Plan for Next Semester (BTP Phase-II)}
\label{sec:roadmap}

The implementation of the Resizer evolution framework is now complete. The planning, coding, and evaluation agents are ready, and the framework can build OpenROAD, test candidate patches, evaluate physical-design results, and preserve the learning from accepted, rejected, and crashed iterations. Therefore, the next phase will focus mainly on running longer evolution experiments and extending the framework beyond the Resizer.

\paragraph{Milestone 1: Large-Scale Resizer Evolution}
I will use the LLM credits to run the Resizer evolution loop for a much larger number of iterations. The objective will be to discover useful changes in gate sizing, buffering, cloning, load splitting, setup repair, and hold repair.

\paragraph{Milestone 2: Complete Graphify Integration}
I will strengthen the integration of Graphify with the agent context. Instead of giving an agent a large collection of unrelated files, Graphify will help identify the functions, callers, data structures, and neighbouring files connected to a proposed change. This source graph will be combined with critical-path reports and previous experiment results. The planning agent will therefore receive both physical evidence from OpenSTA and precise source-code context before selecting a mutation.

\paragraph{Milestone 3: Validation and Integration of Evolved Code}
Promising Resizer candidates will be tested across multiple designs, timing constraints, and technology libraries. The best compatible changes will then be combined in a clean OpenROAD integration branch. This integrated version will be evaluated through the complete OpenROAD flow to check whether local Resizer improvements remain useful after placement, clock-tree synthesis, and routing.

\paragraph{Milestone 4: Expansion to Other OpenROAD Modules}
After establishing stable Resizer results, I will apply the same evolution method to broader OpenROAD directories. Suitable targets include:
\begin{itemize}
  \item \texttt{src/gpl}, for global-placement objectives such as wirelength, density, congestion, and timing;
  \item \texttt{src/cts}, for clock-tree topology, buffer selection, skew, latency, and clock power;
  \item \texttt{src/grt}, for global-routing congestion, overflow, via count, and wirelength;
  \item \texttt{src/drt}, for detailed-routing search, DRC repair, and runtime;
  \item \texttt{src/dpl}, for detailed-placement legalization, displacement, and cell-orientation decisions; and
  \item \texttt{src/mpl}, for macro-placement cost functions and congestion-aware macro movement.
\end{itemize}

These modules will be studied one at a time because each requires a different correctness contract. For example, placement evolution must preserve legality, CTS evolution must control both skew and insertion delay, and routing evolution must not introduce new DRC violations. The Resizer framework will provide the common infrastructure, while module-specific tests, metrics, context, and safety checks will be added through separate adapters.

\paragraph{Milestone 5: A Unified OpenROAD Evolution Framework}
The final goal is to create a common framework capable of evolving several stages of OpenROAD. A high-level planning agent will identify which design stage is limiting the final QoR and route the task to the corresponding module agent. Each module will expose four components: its allowed source files, Graphify context, evaluation metrics, and correctness checks. Accepted module-level changes will be stored as independent champions before being combined and tested in the complete flow.

This design will also make the framework reusable for other EDA tasks. The same loop can be adapted to logic synthesis, RTL optimization, clock-tree construction, placement, routing, DRC repair, power optimization, and design-for-test algorithms. Only the source context, benchmark driver, objective function, and correctness contract need to change. The broader objective of Phase-II is therefore not only to evolve one OpenROAD module, but to develop a reusable method for safely discovering new EDA heuristics through LLM-guided code evolution.


% ==============================================================================
\section{Conclusion}
\label{sec:conclusion}

In this BTP report, I investigated autonomous evolution for OpenROAD physical-design algorithms. The initial \texttt{OptMirror} work improved runtime but produced little WNS/TNS gain. A study of ~3{,}500 isolated patches then showed that most unguided changes regress timing, while a small positive tail justifies continued search. Critical-path analysis connected the recurring slow-slew, weak-driver, deep-logic, and high-fanout failures to the  Resizer, including the 943-fanout Ibex bottleneck. These findings motivated the multi-agent Resizer framework, where every candidate must pass build, regression, placement, electrical, hold, and determinism checks before its timing result can be trusted.

% ==============================================================================
\bibliographystyle{IEEEtran}
\begin{thebibliography}{10}
\footnotesize
\setlength{\itemsep}{0.5pt plus 0.5pt}

\bibitem{selfevolvedabc}
C.~Yu and H.~Ren,
``Autonomous Evolution of EDA Tools: Multi-Agent Self-Evolved ABC,''
\textit{arXiv preprint arXiv:2604.15082}, 2026.

\bibitem{openroad}
A.~B.~Kahng \textit{et al.},
``OpenROAD: Toward a Self-Driving, Open-Source Digital Layout Implementation Tool Chain,''
in \textit{Proc. GOMACTECH}, 2019, pp.~1--5.

\bibitem{alphaevolve}
Google DeepMind,
``AlphaEvolve: A Gemini-Powered Coding Agent for Designing Advanced Algorithms,'' May 2025.

\bibitem{satlution}
C.~Yu, R.~Liang, C.-T.~Ho, and H.~Ren,
``Autonomous Code Evolution Meets NP-Completeness,''
\textit{arXiv preprint arXiv:2509.07367}, 2025.

\bibitem{vprevolve}
Q.~Wu, T.~Jafri, A.~Arora, and V.~A.~Chhabria,
``VPR-Evolve: Multi-Agent Algorithm Evolution for FPGA Place and Route,''
\textit{arXiv preprint arXiv:2607.24998}, 2026.

\bibitem{evodrc}
B.-Y.~Wu, C.-T.~Ho, H.~Yang, B.~Khailany, and V.~A.~Chhabria,
``EvoDRC: Self-Evolving DRC Repair Framework for Advanced Technology Nodes,''
\textit{arXiv preprint arXiv:2607.20019}, 2026.

\bibitem{graphify}
Graphify Labs,
``Graphify: Call-Graph and Semantic Code Indexing for Large C++ Repositories,''
\url{https://github.com/Graphify-Labs/graphify}, 2025.

\end{thebibliography}

\end{document}
